% GENERATED FILE. Edit canonical lessons, then run npm run generate:books.
% canonical-source-hash: fnv1a64:3a99b1538ef86df5
% canonical-lessons: MR-C511-vyavsay, MR-C511-napha, MR-C511-bhandval, MR-C511-bhagidar, MR-C511-guntavne, MR-R511-review-504-511

\chapter{Running a Business}
\label{ch:mr-running-a-business}
\hlchaptermodality{511}

\begin{chapteropening}
\textbf{By the end of this chapter:} \emph{I can talk about a business, its partners and its profit.}

Complete the last lesson of chapter 511: I can talk about a business, its partners and its profit.
\end{chapteropening}

\section[\texorpdfstring{\mr{व्यवसाय} (vyavsāy)}{व्यवसाय (vyavsāy)}]{\mr{व्यवसाय} (vyavsāy) — a business; a profession}

\label{lesson:MR-C511-vyavsay}



\begin{quote}
Before the new one: say the Marathi for to retire, then the Marathi for to dismiss.
\end{quote}

\subsection*{You'll want to know: \mr{व्यवसाय}}
\textbf{\mr{व्यवसाय}} — \emph{vyavsāy} — \textquotedblleft{}a business; a profession\textquotedblright{}.

\begin{quote}
\mr{माझ्या} \mr{काकांचा} \mr{कपड्यांचा} \mr{व्यवसाय} \mr{आहे}.
\end{quote}

\begin{quote}
\emph{mājhyā kākāñcā kapaḍyāñcā vyavsāy āhe.}
\end{quote}

\begin{quote}
My uncle has a clothes business.
\end{quote}

Masculine.

\begin{cousinweb}
\textbf{\mr{व्यवसाय}} (\emph{vyavsāy}) is a tatsama, taken from Sanskrit \emph{vyavasāya}, \textquotedblleft{}effort, undertaking\textquotedblright{} (Monier-Williams).
\end{cousinweb}

\begin{grammarlens}[title={What you've built}]
One more word for describing what you lived and felt.
\end{grammarlens}

\subsection*{Your turn}
\begin{itemize}
\raggedright
  \item \emph{Say it:} \emph{vyavsāy}
  \item \emph{Say it:} \emph{vyavsāy}, once more
  \item \emph{Say it:} \emph{mājhyā kākāñcā kapaḍyāñcā vyavsāy āhe.} — My uncle has a clothes business.
  \item \emph{Recall:} say the Marathi for to retire, then the Marathi for to dismiss, then say \emph{vyavsāy} again
\end{itemize}

\begin{tcolorbox}[breakable,colback=teal!4,colframe=teal!35!black,title={Before you move on}]
What does \mr{व्यवसाय} mean? (\textquotedblleft{}A business; a profession\textquotedblright{}.) Say it once more.
\end{tcolorbox}
\section[\texorpdfstring{\mr{नफा} (naphā)}{नफा (naphā)}]{\mr{नफा} (naphā) — profit}

\label{lesson:MR-C511-napha}



\begin{quote}
Before the new one: say the Marathi for to dismiss, then the Marathi for a business.
\end{quote}

\subsection*{You'll want to know: \mr{नफा}}
\textbf{\mr{नफा}} — \emph{naphā} — \textquotedblleft{}profit\textquotedblright{}.

\begin{quote}
\mr{या} \mr{वर्षी} \mr{दुकानाला} \mr{चांगला} \mr{नफा} \mr{झाला}.
\end{quote}

\begin{quote}
\emph{yā varṣī dukānālā cāṅglā naphā jhālā.}
\end{quote}

\begin{quote}
The shop made a good profit this year.
\end{quote}

Masculine.

Its opposite is \mr{तोटा}: \mr{नफा} \mr{होणे}, \mr{तोटा} \mr{होणे}.

\begin{cousinweb}
\textbf{\mr{नफा}} (\emph{naphā}) is a loan from Persian \emph{nafʿ}, \textquotedblleft{}gain\textquotedblright{}, usually traced to Arabic \emph{nafʿ}, \textquotedblleft{}benefit\textquotedblright{}.
\end{cousinweb}

\begin{grammarlens}[title={What you've built}]
One more word for describing what you lived and felt.
\end{grammarlens}

\subsection*{Your turn}
\begin{itemize}
\raggedright
  \item \emph{Say it:} \emph{naphā}
  \item \emph{Say it:} \emph{naphā}, once more
  \item \emph{Say it:} \emph{yā varṣī dukānālā cāṅglā naphā jhālā.} — The shop made a good profit this year.
  \item \emph{Recall:} say the Marathi for to dismiss, then the Marathi for a business, then say \emph{naphā} again
\end{itemize}

\begin{tcolorbox}[breakable,colback=teal!4,colframe=teal!35!black,title={Before you move on}]
What does \mr{नफा} mean? (\textquotedblleft{}Profit\textquotedblright{}.) Say it once more.
\end{tcolorbox}
\section[\texorpdfstring{\mr{भांडवल} (bhāṇḍval)}{भांडवल (bhāṇḍval)}]{\mr{भांडवल} (bhāṇḍval) — capital; funds for a business}

\label{lesson:MR-C511-bhandval}



\begin{quote}
Before the new one: say the Marathi for a business, then the Marathi for profit.
\end{quote}

\subsection*{You'll want to know: \mr{भांडवल}}
\textbf{\mr{भांडवल}} — \emph{bhāṇḍval} — \textquotedblleft{}capital; funds for a business\textquotedblright{}.

\begin{quote}
\mr{दुकान} \mr{सुरू} \mr{करायला} \mr{त्याला} \mr{भांडवल} \mr{हवं} \mr{होतं}.
\end{quote}

\begin{quote}
\emph{dukān surū karāylā tyālā bhāṇḍval havaṁ hotaṁ.}
\end{quote}

\begin{quote}
He needed capital to start the shop.
\end{quote}

Neuter.

\begin{cousinweb}
\textbf{\mr{भांडवल}} (\emph{bhāṇḍval}) is usually traced to Sanskrit \emph{bhāṇḍa}, \textquotedblleft{}pot, wares, goods\textquotedblright{} (CDIAL 9440), the stock a trader starts with; the ending is Marathi.
\end{cousinweb}

\begin{grammarlens}[title={What you've built}]
One more word for describing what you lived and felt.
\end{grammarlens}

\subsection*{Your turn}
\begin{itemize}
\raggedright
  \item \emph{Say it:} \emph{bhāṇḍval}
  \item \emph{Say it:} \emph{bhāṇḍval}, once more
  \item \emph{Say it:} \emph{dukān surū karāylā tyālā bhāṇḍval havaṁ hotaṁ.} — He needed capital to start the shop.
  \item \emph{Recall:} say the Marathi for a business, then the Marathi for profit, then say \emph{bhāṇḍval} again
\end{itemize}

\begin{tcolorbox}[breakable,colback=teal!4,colframe=teal!35!black,title={Before you move on}]
What does \mr{भांडवल} mean? (\textquotedblleft{}Capital; funds for a business\textquotedblright{}.) Say it once more.
\end{tcolorbox}
\section[\texorpdfstring{\mr{भागीदार} (bhāgīdār)}{भागीदार (bhāgīdār)}]{\mr{भागीदार} (bhāgīdār) — a business partner}

\label{lesson:MR-C511-bhagidar}



\begin{quote}
Before the new one: say the Marathi for profit, then the Marathi for capital.
\end{quote}

\subsection*{You'll want to know: \mr{भागीदार}}
\textbf{\mr{भागीदार}} — \emph{bhāgīdār} — \textquotedblleft{}a business partner\textquotedblright{}.

\begin{quote}
\mr{दोन} \mr{मित्र} \mr{या} \mr{व्यवसायात} \mr{भागीदार} \mr{आहेत}.
\end{quote}

\begin{quote}
\emph{don mitra yā vyavsāyāt bhāgīdār āhet.}
\end{quote}

\begin{quote}
Two friends are partners in this business.
\end{quote}

Masculine.

\begin{cousinweb}
\textbf{\mr{भागीदार}} (\emph{bhāgīdār}) is built on \mr{भाग}, \textquotedblleft{}a share\textquotedblright{}, from Sanskrit \emph{bhāga} (Monier-Williams), with the Persian ending \emph{-dār}, \textquotedblleft{}holder\textquotedblright{}: the one who holds a share.
\end{cousinweb}

\begin{grammarlens}[title={What you've built}]
One more word for describing what you lived and felt.
\end{grammarlens}

\subsection*{Your turn}
\begin{itemize}
\raggedright
  \item \emph{Say it:} \emph{bhāgīdār}
  \item \emph{Say it:} \emph{bhāgīdār}, once more
  \item \emph{Say it:} \emph{don mitra yā vyavsāyāt bhāgīdār āhet.} — Two friends are partners in this business.
  \item \emph{Recall:} say the Marathi for profit, then the Marathi for capital, then say \emph{bhāgīdār} again
\end{itemize}

\begin{tcolorbox}[breakable,colback=teal!4,colframe=teal!35!black,title={Before you move on}]
What does \mr{भागीदार} mean? (\textquotedblleft{}A business partner\textquotedblright{}.) Say it once more.
\end{tcolorbox}
\section[\texorpdfstring{\mr{गुंतवणे} (guntavṇe)}{गुंतवणे (guntavṇe)}]{\mr{गुंतवणे} (guntavṇe) — to invest; to tie up}

\label{lesson:MR-C511-guntavne}



\begin{quote}
Before the new one: say the Marathi for capital, then the Marathi for a business partner.
\end{quote}

\subsection*{You'll want to know: \mr{गुंतवणे}}
\textbf{\mr{गुंतवणे}} — \emph{guntavṇe} — \textquotedblleft{}to invest; to tie up\textquotedblright{}.

\begin{quote}
\mr{भागीदारांनी} \mr{नव्या} \mr{दुकानात} \mr{पैसे} \mr{गुंतवले}.
\end{quote}

\begin{quote}
\emph{bhāgīdārānnī navyā dukānāt paise guntavle.}
\end{quote}

\begin{quote}
The partners invested money in the new shop.
\end{quote}

Past, with -\mr{ने} (he …): \mr{त्याने} \mr{पैसे} \mr{गुंतवले}, \emph{tyāne paise guntavle}.

Where the money goes takes -\mr{त}: \mr{व्यवसायात} \mr{पैसे} \mr{गुंतवणे}.

\begin{cousinweb}
\textbf{\mr{गुंतवणे}} (\emph{guntavṇe}) is built on \mr{गुंतणे}, \textquotedblleft{}to get tied up\textquotedblright{}, as its causative; CDIAL sets \mr{गुंतणे} under Sanskrit \emph{gupháti}, \textquotedblleft{}strings together\textquotedblright{} (CDIAL 4205): invested money is tied up in a business.
\end{cousinweb}

\begin{grammarlens}[title={What you've built}]
One more word for describing what you lived and felt.
\end{grammarlens}

\subsection*{Your turn}
\begin{itemize}
\raggedright
  \item \emph{Say it:} \emph{guntavṇe}
  \item \emph{Say it:} \emph{guntavṇe}, once more
  \item \emph{Say it:} \emph{bhāgīdārānnī navyā dukānāt paise guntavle.} — The partners invested money in the new shop.
  \item \emph{Recall:} say the Marathi for capital, then the Marathi for a business partner, then say \emph{guntavṇe} again
\end{itemize}

\begin{tcolorbox}[breakable,colback=teal!4,colframe=teal!35!black,title={Before you move on}]
What does \mr{गुंतवणे} mean? (\textquotedblleft{}To invest; to tie up\textquotedblright{}.) Say it once more.
\end{tcolorbox}
\section[(dialogue)]{Looking Back at Work and Crafts}

\label{lesson:MR-R511-review-504-511}



\begin{quote}
Say the Marathi for a carpenter and a tailor.
\end{quote}

\subsection*{Your turn}
Say these aloud:

\begin{itemize}
\raggedright
  \item a carpenter, a tailor, a mason, a barber, to hammer in
  \item a potter, a blacksmith, a goldsmith, a weaver, to set up
  \item a clerk, an office attendant, a rubber stamp, a document, to scrutinize
  \item recruitment, employment, training, a knack, to join a post
  \item a wage, daily wages, an honorarium, a pay rise, to earn
  \item a promotion, a position, senior, junior, to strive
  \item time off work, unemployment, to resign, to retire, to dismiss
  \item a business, profit, capital, a business partner, to invest
\end{itemize}

\begin{tcolorbox}[breakable,colback=teal!4,colframe=teal!35!black,title={Before you move on}]
A carpenter? (\textbf{\mr{सुतार}}, \emph{sutār}.) A tailor? (\textbf{\mr{शिंपी}}, \emph{śimpī}.) A mason? (\textbf{\mr{गवंडी}}, \emph{gavaṇḍī}.) A barber? (\textbf{\mr{न्हावी}}, \emph{nhāvī}.) To hammer in? (\textbf{\mr{ठोकणे}}, \emph{ṭhokṇe}.) A potter? (\textbf{\mr{कुंभार}}, \emph{kumbhār}.) A blacksmith? (\textbf{\mr{लोहार}}, \emph{lohār}.) A goldsmith? (\textbf{\mr{सोनार}}, \emph{sonār}.) A weaver? (\textbf{\mr{विणकर}}, \emph{viṇkar}.) To set up? (\textbf{\mr{थाटणे}}, \emph{thāṭṇe}.) A clerk? (\textbf{\mr{लिपिक}}, \emph{lipik}.) An office attendant? (\textbf{\mr{शिपाई}}, \emph{śipāī}.) A rubber stamp? (\textbf{\mr{शिक्का}}, \emph{śikkā}.) A document? (\textbf{\mr{कागदपत्र}}, \emph{kāgadpatra}.) To scrutinize? (\textbf{\mr{छाननी} \mr{करणे}}, \emph{chānnī karṇe}.) Recruitment? (\textbf{\mr{भरती}}, \emph{bhartī}.) Employment? (\textbf{\mr{रोजगार}}, \emph{rojgār}.) Training? (\textbf{\mr{प्रशिक्षण}}, \emph{praśikṣaṇ}.) A knack? (\textbf{\mr{हातोटी}}, \emph{hātoṭī}.) To join a post? (\textbf{\mr{रुजू} \mr{होणे}}, \emph{rujū hoṇe}.) A wage? (\textbf{\mr{वेतन}}, \emph{vetan}.) Daily wages? (\textbf{\mr{मजुरी}}, \emph{majurī}.) An honorarium? (\textbf{\mr{मानधन}}, \emph{māndhan}.) A pay rise? (\textbf{\mr{पगारवाढ}}, \emph{pagārvāḍh}.) To earn? (\textbf{\mr{कमावणे}}, \emph{kamāvṇe}.) A promotion? (\textbf{\mr{पदोन्नती}}, \emph{padonnatī}.) A position? (\textbf{\mr{पद}}, \emph{pad}.) Senior? (\textbf{\mr{वरिष्ठ}}, \emph{variṣṭha}.) Junior? (\textbf{\mr{कनिष्ठ}}, \emph{kaniṣṭha}.) To strive? (\textbf{\mr{झटणे}}, \emph{jhaṭṇe}.) Time off work? (\textbf{\mr{रजा}}, \emph{rajā}.) Unemployment? (\textbf{\mr{बेकारी}}, \emph{bekārī}.) To resign? (\textbf{\mr{राजीनामा} \mr{देणे}}, \emph{rājīnāmā deṇe}.) To retire? (\textbf{\mr{निवृत्त} \mr{होणे}}, \emph{nivṛtta hoṇe}.) To dismiss? (\textbf{\mr{बडतर्फ} \mr{करणे}}, \emph{baḍtarph karṇe}.) A business? (\textbf{\mr{व्यवसाय}}, \emph{vyavsāy}.) Profit? (\textbf{\mr{नफा}}, \emph{naphā}.) Capital? (\textbf{\mr{भांडवल}}, \emph{bhāṇḍval}.) A business partner? (\textbf{\mr{भागीदार}}, \emph{bhāgīdār}.) To invest? (\textbf{\mr{गुंतवणे}}, \emph{guntavṇe}.)
\end{tcolorbox}
